% sec1_intro.tex: title, abstract, Section 1 (Introduction).
% Labels defined here: sec:intro, sec:results, sec:context, sec:strategy, sec:organisation, eq:sigma, eq:jacobi,
%   eq:cp, eq:law, eq:epsLambda, eq:paper1-averaged, def:analytic, thm:analytic, cor:pure, thm:onen, rem:scope,
%   rem:lommel.

\title[Recurrence coefficients with power-law atoms]{Recurrence coefficients of discrete measures with power-law
atoms: pointwise asymptotics and the $1/n$ term}
\author{Gordon [Surname]}
\date{Draft of 15 October 2026}

\begin{abstract}
Let $\sigma=\sum_{k\ge1}k^{-s}\ell(k)\,\delta_{k^{-p}}$ with $p>0$ and $s>1$, so that its Jacobi matrix is a compact
operator with eigenvalues $k^{-p}$. A previous paper established that the recurrence coefficients satisfy
$n^pb_n\to c(p)$ in a double Ces\`aro sense, with an explicit constant $c(p)$ that comes from a constrained equilibrium
problem. Here we prove the pointwise law $n^pb_n=c(p)(1+O(1/n))$ and $a_n/(2b_n)=1+O(1/n)$ for every $p>0$ and $s>1$,
both for pure powers ($\ell\equiv1$) and for every factor $\ell$ that extends to a slowly varying analytic function in a
sector $|\arg t|\le\vartheta$, however small $\vartheta>0$ is. For pure powers with $p>1$ we also find the $1/n$ term:
$n^pb_n=c(p)(1-(s-1)/(2n)+O(n^{-2}))$ and $n^pa_n=2c(p)(1-(s-1+p)/(2n)+O(n^{-2}))$. The proof is a Riemann--Hilbert
analysis of discrete orthogonal polynomials whose nodes accumulate at a hard edge. The outer scalar function depends on
whether $p$ is larger than, equal to or smaller than $1$. At the hard edge a non-analytic slowly varying factor is
handled by a comparison, in the manner of Conway and Deift, with an explicitly solvable Bessel model; this needs the
Bessel order $(s-1-p)/p$ to lie in $(-1,1)$. A Christoffel transformation, which shifts $s$ by $p$ and preserves the
law together with its rate, removes the restriction.
\end{abstract}

\maketitle

\section{Introduction}\label{sec:intro}

A compact Jacobi operator is determined by its spectral measure. If its eigenvalues are $k^{-p}$ and its spectral
weights decay like $k^{-s}$, how do its matrix entries behave? For $p>0$, $s>1$ and a sequence
$\ell\colon\N\to(0,\infty)$ we study the discrete measure
\begin{equation}\label{eq:sigma}
  \sigma_{p,s,\ell}:=\sum_{k\ge1}k^{-s}\ell(k)\,\delta_{k^{-p}},\qquad \sigma_{p,s}:=\sigma_{p,s,1}.
\end{equation}
It lives on $(0,1]$, and its atoms accumulate only at $0$. Let $P_n$ be its monic orthogonal polynomials and
$h_n:=\int P_n^2\diff\sigma$. They satisfy the three-term recurrence
\begin{equation}\label{eq:jacobi}
  xP_n=P_{n+1}+a_nP_n+b_n^2P_{n-1}\qquad(n\ge0,\ b_0:=0,\ b_n>0),
\end{equation}
so that $b_n^2=h_n/h_{n-1}$; the numbers $a_n$ and $b_n$ are the entries of the Jacobi matrix of $\sigma$. With
\begin{equation}\label{eq:cp}
  I(p):=\frac{\sqrt\pi}{2}\,\frac{\Gamma(\frac1p+\frac12)}{\Gamma(\frac1p+1)},\qquad
  c(p):=\frac14\Bigl(\frac{\pi}{2I(p)}\Bigr)^{p},
\end{equation}
for example $c(1)=\frac12$ and $c(2)=\pi^2/16$, the question is whether
\begin{equation}\label{eq:law}
  n^pb_n\longrightarrow c(p)\qquad\text{and}\qquad n^pa_n\longrightarrow2c(p)\qquad(n\to\infty),
\end{equation}
and at what rate. We call \eqref{eq:law} the \emph{pointwise law}. The constant $c(p)$ comes from the companion paper
\cite{Paper1}, which proves the first half of \eqref{eq:law} in an averaged sense (see \eqref{eq:paper1-averaged}
below).

\subsection{Results}\label{sec:results}
The slowly varying factors we allow are those that extend analytically to a sector, written in the logarithmic
variable $\xi=\log t$.

\begin{definition}[analytic slowly varying factors]\label{def:analytic}
Let $\vartheta>0$. A sequence $\ell\colon\N\to(0,\infty)$ is \emph{analytic of angle $\vartheta$} if
$\ell(k)=\exp\Lambda(\log k)$ for all $k\ge1$, where, for some $c_0>0$, the function $\Lambda$ is analytic on
$S_\vartheta:=\{\xi\in\C:\ \re\xi\ge-c_0,\ |\im\xi|\le\vartheta\}$, real on $S_\vartheta\cap\R$, bounded on
$S_\vartheta\cap\{\re\xi\le0\}$, and
\begin{equation}\label{eq:epsLambda}
  \varepsilon_\Lambda(T):=\sup\bigl\{|\Lambda'(\xi)|:\ \xi\in S_\vartheta,\ \re\xi\ge T\bigr\}\longrightarrow0
  \qquad(T\to\infty).
\end{equation}
We then also say that $\Lambda$ is analytic of angle $\vartheta$. If $\vartheta'<\vartheta$, a sequence that is
analytic of angle $\vartheta$ is also analytic of angle $\vartheta'$.
\end{definition}

Such an $\ell$ is slowly varying. The angle may be arbitrarily small, because the Riemann--Hilbert analysis evaluates
the weight only in a sector whose opening we are free to shrink (Lemma~\ref{lem:weightdomain} and
Remark~\ref{rem:angles}). For example, $\exp\sqrt{\log(k+e)}$ and $(\log(k+e))^\beta$, $\beta\in\R$, are analytic of
every angle $\vartheta<\pi$, while $\Lambda(\xi)=\sqrt{\xi+1}+(\xi+2)^{-1/2}\log(1-q\cos2\xi)$ with $q=1/\cosh(0.6)$
has logarithmic branch points at $\xi=\pi m\pm0.3i$, $m\in\mathbb Z$, and is analytic only of the angles
$\vartheta<0.3$ (Section~\ref{sec:numerics}).

\begin{theorem}\label{thm:analytic}
Let $p>0$ and $s>1$, and let $\ell$ be analytic of angle $\vartheta$ for some $\vartheta>0$. Then the Jacobi
coefficients of $\sigma_{p,s,\ell}$ satisfy, as $n\to\infty$,
\[
  n^pb_n=c(p)\bigl(1+O(1/n)\bigr),\qquad \frac{a_n}{2b_n}=1+O(1/n),
\]
with constants depending on $p$, $s$ and $\ell$.
\end{theorem}

The constant sequence $\ell\equiv1$ is analytic of every angle, with $\Lambda\equiv0$.

\begin{corollary}[pure powers]\label{cor:pure}
For every $p>0$ and $s>1$ the measure $\sigma_{p,s}$ satisfies $n^pb_n=c(p)(1+O(1/n))$ and $a_n/(2b_n)=1+O(1/n)$.
In particular, for real $a>0$ and $b$ with $2a+b>1$, the measure $\sigma_{a,2a+b}$ has the moments
$\zeta(aj+2a+b)$, $j\ge0$, and Hankel determinants $\det_{1\le i,j\le n}\zeta(a(i+j)+b)$
\textup{(}\cite[Corollary~1.4 and its proof]{Paper1}\textup{)}, and its recurrence coefficients satisfy
$n^ab_n=c(a)(1+O(1/n))$ and $a_n/(2b_n)=1+O(1/n)$.
\end{corollary}

For pure powers with $p>1$ we identify the next term.

\begin{theorem}\label{thm:onen}
Let $p>1$ and $s>1$. Then the Jacobi coefficients of $\sigma_{p,s}$ satisfy, as $n\to\infty$,
\[
  n^pb_n=c(p)\Bigl(1-\frac{s-1}{2n}+O(n^{-2})\Bigr),\qquad
  n^pa_n=2c(p)\Bigl(1-\frac{s-1+p}{2n}+O(n^{-2})\Bigr).
\]
Consequently $a_n/(2b_n)=1-p/(2n)+O(n^{-2})$ and $b_n=c(p)\bigl(n+\frac{s-1}{2p}\bigr)^{-p}\bigl(1+O(n^{-2})\bigr)$.
\end{theorem}

Here, as everywhere in this paper, $a_n/(2b_n)$ pairs the diagonal entry at $P_n$ with the off-diagonal entry $b_n$
between $P_{n-1}$ and $P_n$, as in \eqref{eq:jacobi}.

The proofs use the following external results, each stated with its hypotheses in
\S\S\ref{sec:tool-wlxz}--\ref{sec:tool-dlmf}. For Theorem~\ref{thm:analytic}: the formulation of the discrete
Riemann--Hilbert problem by Wu, Lin, Xu and Zhao \cite{WLXZ}, as a pattern only (none of their theorems is applied);
the explicit Bessel model of Kuijlaars, McLaughlin, Van Assche and Vanlessen \cite[\S6]{KMVV}, including its freedom
in the choice of angle; Theorem~4.4 of Conway and Deift \cite{CD}, together with the argument of their Theorem~4.7
and Proposition~4.9 and one direction of their Proposition~4.6, which we re-prove; and special-function facts from
\cite{DLMF}. Theorem~\ref{thm:onen} uses in addition the classical left-endpoint Bessel parametrix of
\cite[(6.42)--(6.53)]{KMVV} and the expansions \cite[\S\S9.7, 10.17, 10.40]{DLMF}. The Laurent coefficients of the
first corrections that enter the proof of Theorem~\ref{thm:onen} (Section~\ref{sec:second} and
Appendix~\ref{app:residues}) were computed by exact symbolic series expansion; the coefficient at the hard edge was
also computed by hand.

\begin{remark}[what is, and is not, proved at order $1/n$]\label{rem:scope}
Theorem~\ref{thm:analytic} bounds the deviation from $c(p)$ at order $1/n$ but does not identify the $1/n$ term. The
proof compares the true problem with an exactly solvable comparison problem: relative to it, the non-analytic part of
$\ell$ moves the first two coefficients at infinity by $o(1/n)$ (Section~\ref{sec:hardedge}), but the comparison
problem itself depends on $\ell$ at order $1/n$. So the $O(1/n)$ is absolute, not relative: we do not claim that
$b_n(\sigma_{p,s,\ell})/b_n(\sigma_{p,s})-1=O(\varepsilon_\Lambda(\log n)/n)$. Conjectures for the $1/n$ term when
$p\le1$ or $\ell\not\equiv1$ are formulated in Section~\ref{sec:further} (Conjectures~\ref{conj:onen}
and~\ref{conj:tilt}).
\end{remark}

\begin{remark}[a classical instance]\label{rem:lommel}
For $(p,s)=(2,2)$ the coefficients are explicit. Van Assche \cite[\S3]{VA} records that the orthonormal Lommel
polynomials satisfy $xp_n=\beta_{n+1}p_{n+1}+\beta_np_{n-1}$ with $\beta_{n+1}=[2\sqrt{(n+\nu)(n+\nu+1)}]^{-1}$. For
$\nu=\frac32$ the orthogonality measure $\mu$ is identified as follows. By the Bessel recurrence
\cite[(10.6.1)]{DLMF}, $\mathsf u_n:=\sqrt{n+\nu}\,J_{\nu+n}(1/x)$ solves $x\mathsf u_n=\beta_{n+1}\mathsf
u_{n+1}+\beta_n\mathsf u_{n-1}$ for $n\ge0$ (with $\beta_0:=[2\sqrt{\nu(\nu-1)}]^{-1}$), and it is square summable by
\cite[(10.2.2)]{DLMF}. For $x\notin\R$ the resolvent vector $(\mathcal J-x)^{-1}e_0$ of the Jacobi matrix $\mathcal J$
solves the same recurrence for $n\ge1$ and is square summable, and such solutions form a line, since the Wronskian
$\beta_n(\mathsf u_n\tilde{\mathsf u}_{n-1}-\mathsf u_{n-1}\tilde{\mathsf u}_n)$ of two of them is constant and tends to $0$. Hence
$\int(t-x)^{-1}\diff\mu(t)$, the first entry of the resolvent, equals $-\mathsf u_0/(\beta_0\mathsf u_{-1})=-2\nu
J_\nu(1/x)/J_{\nu-1}(1/x)$. For $\nu=\frac32$ this is $-3(z^{-1}-\cot z)$ with $z=1/x$ \cite[(10.16.1),
(10.6.1)]{DLMF}, and the partial fractions of $\cot$ \cite[(4.22.3)]{DLMF} show that $\mu$ has mass $3/(\pi k)^2$ at
$\pm1/(\pi k)$, $k\ge1$. Its image under $t\mapsto\pi^2t^2$ is $(6/\pi^2)\sigma_{2,2}$. The quadratic transformation
(the monic orthogonal polynomials of $\mu$ of degree $2n$ are $q_n(t^2)$, where $q_n$ are those of the image of $\mu$
under $t\mapsto t^2$) gives $a_n=\pi^2(\beta_{2n}^2+\beta_{2n+1}^2)$ and $b_n=\pi^2\beta_{2n-1}\beta_{2n}$ for
$\sigma_{2,2}$, that is,
\[
  b_n=\frac{\pi^2}{(4n+1)\sqrt{(4n-1)(4n+3)}},\qquad a_n=\frac{2\pi^2}{(4n+1)(4n+5)}\qquad(n\ge1).
\]
Thus $n^2b_n/c(2)=1-\frac1{2n}+O(n^{-2})$ and $a_n/(2b_n)=1-\frac1n+O(n^{-2})$, in agreement with
Theorem~\ref{thm:onen} at $(p,s)=(2,2)$.
\end{remark}

\subsection{Context}\label{sec:context}
\emph{Why the pointwise law is subtle.} By \cite[Corollary~1.2]{Paper1}, specialised to the atoms $x_k=k^{-p}$, if
the weights $w_k\ge0$ satisfy $\sum_kw_k<\infty$ and $w_k\ge c_1k^{-s'}$ for some constants $c_1>0$ and $s'\in\R$,
then the recurrence coefficients of $\sum_kw_k\delta_{k^{-p}}$ satisfy
\begin{equation}\label{eq:paper1-averaged}
  \frac{2}{n^2}\sum_{j=1}^{n-1}(n-j)\log\bigl(j^pb_j\bigr)\longrightarrow\log c(p).
\end{equation}
The constant comes from an explicitly solvable constrained equilibrium problem \cite[Theorem~1.3]{Paper1}, whose
solution $\mu^*$ is the zero distribution that a Jacobi matrix with entries $a_n\approx2b_n\approx2c(p)n^{-p}$ would
produce. But \eqref{eq:paper1-averaged} is a double Ces\`aro mean of $\log(j^pb_j)$, that is, a statement about the
logarithm of the Hankel determinant $D_n$ of $\sigma$ at order $n^2$. Multiplying the weights by factors in
$[e^{-\varepsilon},e^{\varepsilon}]$ changes $\log D_n$ by at most $\varepsilon n$, so the averaged law cannot see such
perturbations, and passing from it to \eqref{eq:law} is a Tauberian step that needs a regularity of $j^pb_j$ which the
averaged argument does not provide; it has to come from the weights. The companion paper \cite{Paper3}
uses the results of this paper to treat general slowly varying and rough weights, and shows there that bounded rough
weights can break the pointwise law.

\emph{Discrete orthogonal polynomials.} The Riemann--Hilbert approach to discrete orthogonal polynomials was developed
by Baik, Kriecherbauer, McLaughlin and Miller \cite{BKMM} for measures on finite node sets in a bounded interval, with
real-analytic node density and potential; their analysis includes saturated regions, which also appear here. Within
the literature we have read, the only Riemann--Hilbert analysis of discrete orthogonal polynomials on infinitely many
nodes accumulating at a finite point is that of Wu, Lin, Xu and Zhao \cite{WLXZ}, for the Tricomi--Carlitz
polynomials; our frame follows theirs. Dai and Yao make the same observation: ``Wu et al.\ [21] studied the case where
the orthogonality lattice consists of infinite nodes with an accumulation point. Except [21], the orthogonality
lattices are always equally spaced'' \cite[\S1]{DY}, [21] being \cite{WLXZ}.

\emph{Non-analytic weights.} Our weights are not analytic at the hard edge, and the standard local parametrix does not
exist there. Conway and Deift \cite{CD} met the same difficulty for the weight $\log\frac{2k}{1-x}$ on $(-1,1)$, which
carries a slowly varying factor at the endpoint $x=1$; they write that ``there is no known parametrix for the
Riemann--Hilbert problem in a neighborhood of the logarithmic singularity at $x=1$'', and they compare the problem with
the Legendre problem through the operator theory behind Riemann--Hilbert problems instead of through a quotient. The
comparison argument of Section~\ref{sec:hardedge} is theirs \cite[Theorem~4.7 and Proposition~4.9]{CD}. What we add
is a comparison model built from the Szeg\H{o} function of the full weight (an explicit parametrix rather than a second
orthogonality problem), a uniform resolvent bound for a Bessel edge of any order $|\alpha|<1$ (rather than the Legendre
value $\alpha=0$), and the treatment of an arbitrary analytic slowly varying factor (rather than one explicit
logarithm). Deift and Piorkowski \cite{DP} extended the
analysis of \cite{CD} to the weight $\log\frac{2}{1-x}$ and write, about a comparison with a weight that has a simple
zero at the endpoint: ``the crucial Theorem~4.7 in [4] requires a different proof in this case: now one needs to
control the behaviour of the Cauchy operator acting on spaces with Muckenhoupt weights'' \cite[\S1]{DP}, [4] being
\cite{CD}. This is consistent with our restriction to Bessel orders $|\alpha|<1$ in the comparison step
(Remark~\ref{rem:alpha-ge-1}). Other approaches to non-analytic weights are Piorkowski's parametrix-free method, which
needs an a priori estimate for the exact solution \cite{Pio}; the $\bar\partial$ methods of McLaughlin and Miller
\cite{MM} and of Trogdon \cite{Tro}, which assume a finite number of Lipschitz derivatives of the weight (for
\cite{Tro}, $C^{k,1}$ regularity up to the endpoints of the support); and the approximation of a smooth symbol by
analytic ones in Deift, Its and Krasovsky \cite[\S6.2]{DIK}.

\emph{Recurrence asymptotics without analyticity.} For exponential weights $W=e^{-Q}$, Levin and Lubinsky \cite{LL}
prove $o(1)$ asymptotics of the recurrence coefficients, with constants given by equilibrium quantities, under
regularity conditions that do not involve analyticity, and without Riemann--Hilbert methods.

\emph{Special cases.} Besides the Lommel case (Remark~\ref{rem:lommel}), we mention the Hankel determinants of zeta
values studied by Monien \cite{Monien} and by Haynes and Zudilin \cite{HaynesZudilin}; for $a=1$ Monien reports a full
asymptotic expansion of the determinants, obtained numerically by D.~Zagier, which goes beyond
Corollary~\ref{cor:pure} but is not proved. The Tricomi--Carlitz polynomials are orthogonal with respect to a measure on the points $\pm1/\sqrt{k+t_0}$, for a
parameter $t_0>0$ \cite[\S3]{VA}; in the squared variable this is an instance of $p=1$ with weights that are not pure
powers. Their recurrence coefficients are explicit \cite[\S3]{VA}, and \cite{WLXZ} obtain uniform asymptotics of the
polynomials.

\subsection{Method}\label{sec:strategy}
The proof of Theorem~\ref{thm:analytic} has an elementary part and a Riemann--Hilbert part. The elementary part is the
Christoffel shift, Theorem~\ref{thm:shift}: the pointwise law passes from $\sigma$ to $x\sigma$ together with any
regular rate, and $x\sigma_{p,s,\ell}=\sigma_{p,s+p,\ell}$. What is new there is the rate; the statement without rate
follows from a proposition of Dai, Piao and Qiao \cite[Proposition~2.1]{DPQ} and the classical transfer formulas
\cite[\S1]{DD}. The shift reduces Theorem~\ref{thm:analytic} to the window $1<s<1+2p$, where the Bessel order
$\alpha:=(s-1-p)/p$ of the hard edge lies in $(-1,1)$ (Theorem~\ref{thm:window}). For the Riemann--Hilbert part we
rescale $z\mapsto n^pz$. The equilibrium measure $\mu^*$ of \cite{Paper1} then has the band $(0,4c(p))$, where the
polynomials oscillate, the saturated region $(4c(p),\infty)$, where their zeros follow the nodes, and a soft edge at
$4c(p)$, and on it we build one frame for all $p>0$ (Section~\ref{sec:frame}). The three geometries $p>1$, $p=1$ and
$p<1$ enter only through the outer scalar function that removes the Gamma-function growth of the node masses
(Section~\ref{sec:outer} and Appendix~\ref{app:outer}): the index function $t(z)=nz^{-1/p}$ maps the slit plane onto a
sector that lies inside, exactly fills, or exceeds the sector where Stirling's formula holds. At the hard edge the
factor enters through $r_n(z)=\ell(nz^{-1/p})/\ell(n)$, which is not analytic at $0$, so the classical Bessel
parametrix is not available. Following Conway and Deift, we compare the problem with a Bessel model built from the
Szeg\H{o} function of the whole band weight; the two problems differ on the lips of the lens by $o(1)$ without a rate,
and for $|\alpha|<1$ a uniform resolvent bound and Bessel localisation turn this into an $O(1/n)$ error in the
coefficients (Section~\ref{sec:hardedge}). For $\ell\equiv1$ and $p>1$ the classical Bessel parametrix does apply, for
every $\alpha>-1$ (Section~\ref{sec:classical}), and it gives the $1/n$ term (Section~\ref{sec:second}): the band factor
is constant up to $O(n^{-p})$, the hard-edge correction is annihilated by both functionals that extract $a_n$ and $b_n$,
and the $1/n$ term comes from the Airy correction at the soft edge. This module cannot handle a non-analytic $r_n$, it
fails at $p=1$, and we do not use it for $p<1$, where the band factor is not constant to $o(1/n)$ near $0$.

\subsection{Organisation}\label{sec:organisation}
Section~\ref{sec:prelim} fixes conventions and states the inputs, and Section~\ref{sec:shift} proves
Theorem~\ref{thm:shift}. Sections~\ref{sec:frame}--\ref{sec:extraction} prove Theorem~\ref{thm:window} and, with the
shift, Theorem~\ref{thm:analytic} and Corollary~\ref{cor:pure}; the interface between the frame and the hard edge is
Proposition~\ref{prop:frame}. Sections~\ref{sec:classical} and~\ref{sec:second} prove Theorem~\ref{thm:onen},
Section~\ref{sec:numerics} gives numerical illustrations, and Section~\ref{sec:further} contains conjectures and open
problems. A reader interested only in Theorem~\ref{thm:onen} needs Sections~\ref{sec:prelim}, \ref{sec:frame},
\ref{sec:outer}, \ref{sec:classical} and~\ref{sec:second} and Appendices~\ref{app:outer} and~\ref{app:residues}, and
nothing from Section~\ref{sec:hardedge}.

Throughout, $n$ is the degree, and $C,C'$ denote positive constants that may change from line to line. They are
never written $c$, which is reserved for $c(p)$; the exceptions are $c_0$ in Definition~\ref{def:analytic} and $c_1$ in
\eqref{eq:paper1-averaged}. Our atoms are exactly $k^{-p}$, so the atom constant of \cite[Theorem~1.1]{Paper1} does not
appear.
